\MakeAbstract{
    在我国高校普遍的多校区办学模式下，校区间通勤的需求不断增加，班车、公共交通与临时拼车并存，但相关信息长期分散在不同渠道中，师生在出行前往往需要自行完成信息检索、路线比较和同行沟通，使用成本较高。针对这一实际需求，本文设计并实现了一个基于 Spring Boot 的校内综合交通平台。平台采用前后端分离架构，以微信小程序作为统一用户入口，后端按用户、拼车、论坛、聊天、文件和交通等业务划分服务，并结合 MySQL 与 MongoDB 构建混合数据存储方案。

    在功能实现上，平台以拼车服务为核心，围绕校园场景下常见的跨校区通勤需求，构建了订单发布、候选查询、申请审批、行程完成和反馈更新的完整业务流程。针对拼车匹配中既存在明确时空约束、又存在用户偏好差异的特点，系统设计了“规则筛选 + 随机森林排序”的二层匹配机制：首先根据路线和时间约束对订单进行可行性筛选；随后结合空间匹配、时间重叠和用户信用等特征，对候选订单进行排序。同时，平台还实现了论坛交流、私聊沟通、交通时刻表查询和行程反馈等功能，以支撑校园出行中的信息共享与协同确认。

    测试结果表明，系统能够较稳定地完成用户登录、拼车管理、论坛互动、即时通信和交通信息查询等主要业务流程，拼车匹配结果与场景约束基本一致，关键查询接口在给定并发条件下保持了可接受的响应性能。本文工作的意义在于面向校园综合交通这一具体应用场景，完成了多源出行信息整合、拼车流程组织及候选排序方法的工程实现，可为高校内部出行服务平台建设提供参考。
}{校内综合交通平台，微信小程序，Spring Boot，拼车匹配，随机森林}

\MakeAbstractEng{
    Under the commonly adopted multi-campus operation model in Chinese universities, the demand for commuting between campuses is continuously increasing, with shuttle buses, public transportation, and ad-hoc carpooling coexisting. However, relevant information has long been scattered across different channels, and faculty and students often need to search for information, compare routes, and coordinate with fellow travelers on their own before traveling, resulting in high usage costs. To address this practical need, this paper designs and implements a campus integrated transportation platform based on Spring Boot. The platform adopts a front-end/back-end separated architecture, with a WeChat Mini Program as the unified user interface. The back end organizes services by business domain, including users, carpooling, forums, chat, files, and transportation, and employs a hybrid data storage solution combining MySQL and MongoDB.
    
    In terms of functional implementation, the platform takes the carpool service as its core and builds a complete business process covering order publishing, candidate query, application approval, trip completion, and feedback update for common inter-campus commuting needs. To address the characteristics of carpool matching, where explicit spatiotemporal constraints coexist with differences in user preferences, the system employs a two-layer matching mechanism combining rule-based filtering and random-forest ranking. Orders are first filtered according to route and time constraints, and feasible candidates are then ranked by features such as spatial matching, time overlap, and user credit. In addition, the platform provides forum interaction, private messaging, timetable query, and trip feedback to support information sharing and collaborative confirmation in campus travel.

    Experimental results show that the platform can stably support major business processes, including user login, carpool management, forum interaction, instant messaging, and transport information query. The carpool matching results remain basically consistent with scenario constraints, and key query interfaces maintain acceptable response performance under the given concurrency setting. The significance of this work lies in completing an engineering implementation for multi-source travel information integration, carpool workflow organization, and candidate ranking in the concrete application scenario of campus comprehensive transportation, which can provide a reference for the construction of travel service platforms in universities.
}{Campus Comprehensive Transportation Platform, WeChat Mini Program, Spring Boot, Carpool Matching, Random Forest}
